%MakeMMrefs.tex

\begin{lemma}\otherTitle{3anass} Associative law for triple conjunction. 
(Contributed by NM,
     8-Apr-1994.)
\begin{align}
\vdash ( ( {\color{formulacolor}\varphi} \wedge {\color{formulacolor}\psi}
    \wedge {\color{formulacolor}\chi} ) \leftrightarrow (
    {\color{formulacolor}\varphi} \wedge ( {\color{formulacolor}\psi} \wedge
    {\color{formulacolor}\chi} ) ) )\label{eq:3anass}\tag{3anass}
\end{align}
\end{lemma}


\begin{lemma}\otherTitle{3orass} Associative law for triple disjunction. 
(Contributed by NM,
     8-Apr-1994.)
\begin{align}
\vdash ( ( {\color{formulacolor}\varphi} \vee {\color{formulacolor}\psi} \vee
    {\color{formulacolor}\chi} ) \leftrightarrow (
    {\color{formulacolor}\varphi} \vee ( {\color{formulacolor}\psi} \vee
    {\color{formulacolor}\chi} ) ) )\label{eq:3orass}\tag{3orass}
\end{align}
\end{lemma}


\begin{lemma}\otherTitle{3bitri} A chained inference from transitive law for
logical equivalence.
       (Contributed by NM, 5-Aug-1993.)
\begin{align}
\vdash ( {\color{formulacolor}\varphi} \leftrightarrow
    {\color{formulacolor}\psi} )\quad\&\quad \vdash (
    {\color{formulacolor}\psi} \leftrightarrow {\color{formulacolor}\chi}
    )\quad\&\quad \vdash ( {\color{formulacolor}\chi} \leftrightarrow
    {\color{formulacolor}\theta} )\quad\Rightarrow\quad \vdash (
    {\color{formulacolor}\varphi} \leftrightarrow {\color{formulacolor}\theta}
    )\label{eq:3bitri}\tag{3bitri}
\end{align}
\end{lemma}


\begin{lemma}\otherTitle{3imtr3i} A mixed syllogism inference, useful for
removing a definition from both
       sides of an implication.  (Contributed by NM, 10-Aug-1994.)
\begin{align}
\vdash ( {\color{formulacolor}\varphi} \rightarrow {\color{formulacolor}\psi}
    )\quad\&\quad \vdash ( {\color{formulacolor}\varphi} \leftrightarrow
    {\color{formulacolor}\chi} )\quad\&\quad \vdash (
    {\color{formulacolor}\psi} \leftrightarrow {\color{formulacolor}\theta}
    )\quad\Rightarrow\quad \vdash ( {\color{formulacolor}\chi} \rightarrow
    {\color{formulacolor}\theta} )\label{eq:3imtr3i}\tag{3imtr3i}
\end{align}
\end{lemma}


\begin{lemma}\otherTitle{3imtr4i} A mixed syllogism inference, useful for
applying a definition to both
       sides of an implication.  (Contributed by NM, 5-Aug-1993.)
\begin{align}
\vdash ( {\color{formulacolor}\varphi} \rightarrow {\color{formulacolor}\psi}
    )\quad\&\quad \vdash ( {\color{formulacolor}\chi} \leftrightarrow
    {\color{formulacolor}\varphi} )\quad\&\quad \vdash (
    {\color{formulacolor}\theta} \leftrightarrow {\color{formulacolor}\psi}
    )\quad\Rightarrow\quad \vdash ( {\color{formulacolor}\chi} \rightarrow
    {\color{formulacolor}\theta} )\label{eq:3imtr4i}\tag{3imtr4i}
\end{align}
\end{lemma}


\begin{lemma}\otherTitle{a1bi} Inference rule introducing a theorem as an
antecedent.  (Contributed by
       NM, 5-Aug-1993.)  (Proof shortened by Wolf Lammen, 11-Nov-2012.)
\begin{align}
\vdash {\color{formulacolor}\varphi}\quad\Rightarrow\quad \vdash (
    {\color{formulacolor}\psi} \leftrightarrow ( {\color{formulacolor}\varphi}
    \rightarrow {\color{formulacolor}\psi} ) )\label{eq:a1bi}\tag{a1bi}
\end{align}
\end{lemma}


\begin{lemma}\otherTitle{a1tru} Anything implies $\top$.  (Contributed by FL,
20-Mar-2011.)  (Proof
     shortened by Anthony Hart, 1-Aug-2011.)
\begin{align}
\vdash ( {\color{formulacolor}\varphi} \rightarrow \top )\label{eq:a1tru}\tag{a1tru}
\end{align}
\end{lemma}


\begin{lemma}\otherTitle{anbi12i} Conjoin both sides of two equivalences. 
(Contributed by NM,
       5-Aug-1993.)
\begin{align}
\vdash ( {\color{formulacolor}\varphi} \leftrightarrow
    {\color{formulacolor}\psi} )\quad\&\quad \vdash (
    {\color{formulacolor}\chi} \leftrightarrow {\color{formulacolor}\theta}
    )\quad\Rightarrow\quad \vdash ( ( {\color{formulacolor}\varphi} \wedge
    {\color{formulacolor}\chi} ) \leftrightarrow ( {\color{formulacolor}\psi}
    \wedge {\color{formulacolor}\theta} ) )\label{eq:anbi12i}\tag{anbi12i}
\end{align}
\end{lemma}


\begin{lemma}\otherTitle{anbi1i} Introduce a right conjunct to both sides of a
logical equivalence.
       (Contributed by NM, 5-Aug-1993.)  (Proof shortened by Wolf Lammen,
       16-Nov-2013.)
\begin{align}
\vdash ( {\color{formulacolor}\varphi} \leftrightarrow
    {\color{formulacolor}\psi} )\quad\Rightarrow\quad \vdash ( (
    {\color{formulacolor}\varphi} \wedge {\color{formulacolor}\chi} )
    \leftrightarrow ( {\color{formulacolor}\psi} \wedge
    {\color{formulacolor}\chi} ) )\label{eq:anbi1i}\tag{anbi1i}
\end{align}
\end{lemma}


\begin{lemma}\otherTitle{anbi2i} Introduce a left conjunct to both sides of a
logical equivalence.
       (Contributed by NM, 5-Aug-1993.)  (Proof shortened by Wolf Lammen,
       16-Nov-2013.)
\begin{align}
\vdash ( {\color{formulacolor}\varphi} \leftrightarrow
    {\color{formulacolor}\psi} )\quad\Rightarrow\quad \vdash ( (
    {\color{formulacolor}\chi} \wedge {\color{formulacolor}\varphi} )
    \leftrightarrow ( {\color{formulacolor}\chi} \wedge
    {\color{formulacolor}\psi} ) )\label{eq:anbi2i}\tag{anbi2i}
\end{align}
\end{lemma}


\begin{lemma}\otherTitle{ancom} Commutative law for conjunction.  Theorem *4.3
of [WhiteheadRussell]
     p. 118.  (Contributed by NM, 25-Jun-1998.)  (Proof shortened by Wolf
     Lammen, 4-Nov-2012.)
\begin{align}
\vdash ( ( {\color{formulacolor}\varphi} \wedge {\color{formulacolor}\psi} )
    \leftrightarrow ( {\color{formulacolor}\psi} \wedge
    {\color{formulacolor}\varphi} ) )\label{eq:ancom}\tag{ancom}
\end{align}
\end{lemma}


\begin{lemma}\otherTitle{andir} Distributive law for conjunction.  (Contributed
by NM, 12-Aug-1994.)
\begin{align}
\vdash ( ( ( {\color{formulacolor}\varphi} \vee {\color{formulacolor}\psi} )
    \wedge {\color{formulacolor}\chi} ) \leftrightarrow ( (
    {\color{formulacolor}\varphi} \wedge {\color{formulacolor}\chi} ) \vee (
    {\color{formulacolor}\psi} \wedge {\color{formulacolor}\chi} ) )
    )\label{eq:andir}\tag{andir}
\end{align}
\end{lemma}


\begin{lemma}\otherTitle{anidm} Idempotent law for conjunction.  (Contributed
by NM, 5-Aug-1993.)  (Proof
     shortened by Wolf Lammen, 14-Mar-2014.)
\begin{align}
\vdash ( ( {\color{formulacolor}\varphi} \wedge {\color{formulacolor}\varphi} )
    \leftrightarrow {\color{formulacolor}\varphi} )\label{eq:anidm}\tag{anidm}
\end{align}
\end{lemma}


\begin{lemma}\otherTitle{ax-1} Axiom \{\em Simp\}.  Axiom A1 of [Margaris] p.
49.  One of the 3 axioms of
     propositional calculus.  The 3 axioms are also given as Definition 2.1 of
     [Hamilton] p. 28.  This axiom is called \{\em Simp\} or "the principle of
     simplification" in \{\em Principia Mathematica\} (Theorem *2.02 of
     [WhiteheadRussell] p. 100) because "it enables us to pass from the joint
     assertion of ${\color{formulacolor}\varphi}$ and
${\color{formulacolor}\psi}$ to the assertion of
${\color{formulacolor}\varphi}$ simply."
     (Contributed by NM, 5-Aug-1993.)
\begin{align}
\vdash ( {\color{formulacolor}\varphi} \rightarrow ( {\color{formulacolor}\psi}
    \rightarrow {\color{formulacolor}\varphi} ) )\label{eq:ax-1}\tag{ax-1}
\end{align}
\end{lemma}


\begin{lemma}\otherTitle{ax-mp} Rule of Modus Ponens.  The postulated inference
rule of propositional
       calculus.  See e.g.  Rule 1 of [Hamilton] p. 73.  The rule says, "if
       ${\color{formulacolor}\varphi}$ is true, and
${\color{formulacolor}\varphi}$ implies ${\color{formulacolor}\psi}$, then
${\color{formulacolor}\psi}$ must also be
       true."  This rule is sometimes called "detachment," since it detaches
       the minor premise from the major premise.  "Modus ponens" is short for
       "modus ponendo ponens," a Latin phrase that means "the mood that by
       affirming affirms" [Sanford] p. 39.  This rule is similar to the rule
of
       modus tollens {\tt mto}.

       Note:  In some web page displays such as the Statement List, the
symbols
       "\&" and "=\}" informally indicate the relationship between the
hypotheses
       and the assertion (conclusion), abbreviating the English words "and"
and
       "implies."  They are not part of the formal language.  (Contributed by
       NM, 5-Aug-1993.)
\begin{align}
\vdash {\color{formulacolor}\varphi}\quad\&\quad \vdash (
    {\color{formulacolor}\varphi} \rightarrow {\color{formulacolor}\psi}
    )\quad\Rightarrow\quad \vdash {\color{formulacolor}\psi}\label{eq:ax-mp}\tag{ax-mp}
\end{align}
\end{lemma}


\begin{lemma}\otherTitle{bicomi} Inference from commutative law for logical
equivalence.  (Contributed by
       NM, 5-Aug-1993.)
\begin{align}
\vdash ( {\color{formulacolor}\varphi} \leftrightarrow
    {\color{formulacolor}\psi} )\quad\Rightarrow\quad \vdash (
    {\color{formulacolor}\psi} \leftrightarrow {\color{formulacolor}\varphi}
    )\label{eq:bicomi}\tag{bicomi}
\end{align}
\end{lemma}


\begin{lemma}\otherTitle{biid} Principle of identity for logical equivalence. 
Theorem *4.2 of
     [WhiteheadRussell] p. 117.  (Contributed by NM, 5-Aug-1993.)
\begin{align}
\vdash ( {\color{formulacolor}\varphi} \leftrightarrow
    {\color{formulacolor}\varphi} )\label{eq:biid}\tag{biid}
\end{align}
\end{lemma}


\begin{lemma}\otherTitle{biimpi} Infer an implication from a logical
equivalence.  (Contributed by NM,
       5-Aug-1993.)
\begin{align}
\vdash ( {\color{formulacolor}\varphi} \leftrightarrow
    {\color{formulacolor}\psi} )\quad\Rightarrow\quad \vdash (
    {\color{formulacolor}\varphi} \rightarrow {\color{formulacolor}\psi}
    )\label{eq:biimpi}\tag{biimpi}
\end{align}
\end{lemma}


\begin{lemma}\otherTitle{biimpri} Infer a converse implication from a logical
equivalence.  (Contributed
       by NM, 5-Aug-1993.)  (Proof shortened by Wolf Lammen, 16-Sep-2013.)
\begin{align}
\vdash ( {\color{formulacolor}\varphi} \leftrightarrow
    {\color{formulacolor}\psi} )\quad\Rightarrow\quad \vdash (
    {\color{formulacolor}\psi} \rightarrow {\color{formulacolor}\varphi}
    )\label{eq:biimpri}\tag{biimpri}
\end{align}
\end{lemma}


\begin{lemma}\otherTitle{bitr2i} An inference from transitive law for logical
equivalence.  (Contributed
       by NM, 5-Aug-1993.)
\begin{align}
\vdash ( {\color{formulacolor}\varphi} \leftrightarrow
    {\color{formulacolor}\psi} )\quad\&\quad \vdash (
    {\color{formulacolor}\psi} \leftrightarrow {\color{formulacolor}\chi}
    )\quad\Rightarrow\quad \vdash ( {\color{formulacolor}\chi} \leftrightarrow
    {\color{formulacolor}\varphi} )\label{eq:bitr2i}\tag{bitr2i}
\end{align}
\end{lemma}


\begin{lemma}\otherTitle{bitr3} Closed nested implication form of {\tt bitr3i}.
Derived automatically from
     {\tt bitr3VD}.  (Contributed by Alan Sare, 31-Dec-2011.)
     (Proof modification is discouraged.)  (New usage is discouraged.)
\begin{align}
\vdash ( ( {\color{formulacolor}\varphi} \leftrightarrow
    {\color{formulacolor}\psi} ) \rightarrow ( ( {\color{formulacolor}\varphi}
    \leftrightarrow {\color{formulacolor}\chi} ) \rightarrow (
    {\color{formulacolor}\psi} \leftrightarrow {\color{formulacolor}\chi} ) )
    )\label{eq:bitr3}\tag{bitr3}
\end{align}
\end{lemma}


\begin{lemma}\otherTitle{bitr3i} An inference from transitive law for logical
equivalence.  (Contributed
       by NM, 5-Aug-1993.)
\begin{align}
\vdash ( {\color{formulacolor}\psi} \leftrightarrow
    {\color{formulacolor}\varphi} )\quad\&\quad \vdash (
    {\color{formulacolor}\psi} \leftrightarrow {\color{formulacolor}\chi}
    )\quad\Rightarrow\quad \vdash ( {\color{formulacolor}\varphi}
    \leftrightarrow {\color{formulacolor}\chi} )\label{eq:bitr3i}\tag{bitr3i}
\end{align}
\end{lemma}


\begin{lemma}\otherTitle{bitr4i} An inference from transitive law for logical
equivalence.  (Contributed
       by NM, 5-Aug-1993.)
\begin{align}
\vdash ( {\color{formulacolor}\varphi} \leftrightarrow
    {\color{formulacolor}\psi} )\quad\&\quad \vdash (
    {\color{formulacolor}\chi} \leftrightarrow {\color{formulacolor}\psi}
    )\quad\Rightarrow\quad \vdash ( {\color{formulacolor}\varphi}
    \leftrightarrow {\color{formulacolor}\chi} )\label{eq:bitr4i}\tag{bitr4i}
\end{align}
\end{lemma}


\begin{lemma}\otherTitle{bitri} An inference from transitive law for logical
equivalence.  (Contributed
       by NM, 5-Aug-1993.)  (Proof shortened by Wolf Lammen, 13-Oct-2012.)
\begin{align}
\vdash ( {\color{formulacolor}\varphi} \leftrightarrow
    {\color{formulacolor}\psi} )\quad\&\quad \vdash (
    {\color{formulacolor}\psi} \leftrightarrow {\color{formulacolor}\chi}
    )\quad\Rightarrow\quad \vdash ( {\color{formulacolor}\varphi}
    \leftrightarrow {\color{formulacolor}\chi} )\label{eq:bitri}\tag{bitri}
\end{align}
\end{lemma}


\begin{lemma}\otherTitle{con1bii} A contraposition inference.  (Contributed by
NM, 5-Aug-1993.)  (Proof
       shortened by Wolf Lammen, 13-Oct-2012.)
\begin{align}
\vdash ( \lnot {\color{formulacolor}\varphi} \leftrightarrow
    {\color{formulacolor}\psi} )\quad\Rightarrow\quad \vdash ( \lnot
    {\color{formulacolor}\psi} \leftrightarrow {\color{formulacolor}\varphi}
    )\label{eq:con1bii}\tag{con1bii}
\end{align}
\end{lemma}


\begin{lemma}\otherTitle{con4bii} A contraposition inference.  (Contributed by
NM, 21-May-1994.)
\begin{align}
\vdash ( \lnot {\color{formulacolor}\varphi} \leftrightarrow \lnot
    {\color{formulacolor}\psi} )\quad\Rightarrow\quad \vdash (
    {\color{formulacolor}\varphi} \leftrightarrow {\color{formulacolor}\psi}
    )\label{eq:con4bii}\tag{con4bii}
\end{align}
\end{lemma}


\begin{lemma}\otherTitle{df-an} Define conjunction (logical 'and').  Definition
of [Margaris] p. 49.  When
     both the left and right operand are true, the result is true; when either
     is false, the result is false.  For example, it is true that
     $( 2 = 2 \wedge 3 = 3 )$.  After we define the constant true $\top$
     ({\tt df-tru}) and the constant false $\bot$ ({\tt df-fal}), we will be
able
     to prove these truth table values: $( ( \top \wedge \top )
\leftrightarrow \top )$
     ({\tt truantru}), $( ( \top \wedge \bot ) \leftrightarrow \bot )$ ({\tt
truanfal}),
     $( ( \bot \wedge \top ) \leftrightarrow \bot )$ ({\tt falantru}), and
     $( ( \bot \wedge \bot ) \leftrightarrow \bot )$ ({\tt falanfal}).

     Contrast with $\vee$ ({\tt df-or}), $\rightarrow$ ({\tt wi}), $\barwedge$
({\tt df-nan}),
     and $\veebar$ ({\tt df-xor}) .  (Contributed by NM, 5-Aug-1993.)
\begin{align}
\vdash ( ( {\color{formulacolor}\varphi} \wedge {\color{formulacolor}\psi} )
    \leftrightarrow \lnot ( {\color{formulacolor}\varphi} \rightarrow \lnot
    {\color{formulacolor}\psi} ) )\label{eq:df-an}
\end{align}
\end{lemma}


\begin{lemma}\otherTitle{df-or} Define disjunction (logical 'or').  Definition
of [Margaris] p. 49.  When
     the left operand, right operand, or both are true, the result is true;
     when both sides are false, the result is false.  For example, it is true
     that $( 2 = 3 \vee 4 = 4 )$ ({\tt ex-or}).  After we define the constant
     true $\top$ ({\tt df-tru}) and the constant false $\bot$ ({\tt df-fal}),
we
     will be able to prove these truth table values:
     $( ( \top \vee \top ) \leftrightarrow \top )$ ({\tt truortru}), $( ( \top
\vee \bot ) \leftrightarrow \top )$
     ({\tt truorfal}), $( ( \bot \vee \top ) \leftrightarrow \top )$ ({\tt
falortru}), and
     $( ( \bot \vee \bot ) \leftrightarrow \bot )$ ({\tt falorfal}).

     This is our first use of the biconditional connective in a definition; we
     use the biconditional connective in place of the traditional "\{=def=\}",
     which means the same thing, except that we can manipulate the
     biconditional connective directly in proofs rather than having to rely on
     an informal definition substitution rule.  Note that if we mechanically
     substitute $( \lnot {\color{formulacolor}\varphi} \rightarrow
{\color{formulacolor}\psi} )$ for $( {\color{formulacolor}\varphi} \vee
{\color{formulacolor}\psi} )$, we end up with an
     instance of previously proved theorem {\tt biid}.  This is the
justification
     for the definition, along with the fact that it introduces a new symbol
     $\vee$.  Contrast with $\wedge$ ({\tt df-an}), $\rightarrow$ ({\tt wi}),
$\barwedge$
     ({\tt df-nan}), and $\veebar$ ({\tt df-xor}) .  (Contributed by NM,
     5-Aug-1993.)
\begin{align}
\vdash ( ( {\color{formulacolor}\varphi} \vee {\color{formulacolor}\psi} )
    \leftrightarrow ( \lnot {\color{formulacolor}\varphi} \rightarrow
    {\color{formulacolor}\psi} ) )\label{df-or}
\end{align}
\end{lemma}


\begin{lemma}\otherTitle{df-tru} Definition of $\top$, a tautology. $\top$ is a
constant true.  In this
     definition {\tt biid} is used as an antecedent, however, any true wff,
such as
     an axiom, can be used in its place.  (Contributed by Anthony Hart,
     13-Oct-2010.)
\begin{align}
\vdash ( \top \leftrightarrow ( {\color{formulacolor}\varphi} \leftrightarrow
    {\color{formulacolor}\varphi} ) )\label{df-tru}
\end{align}
\end{lemma}


\begin{lemma}\otherTitle{eqid} Law of identity (reflexivity of class equality).
Theorem 6.4 of [Quine]
       p. 41.

       This law is thought to have originated with Aristotle (\{\em
Metaphysics\},
       Book VII, Part 17).  (Thanks to Stefan Allan for this information.)
       (Contributed by NM, 5-Aug-1993.)
\begin{align}
\vdash {\color{classcolor}A} = {\color{classcolor}A}\label{eq:eqid}\tag{eqid}
\end{align}
\end{lemma}


\begin{lemma}\otherTitle{exmid} Law of excluded middle, also called the
principle of \{\em tertium non datur\}.
     Theorem *2.11 of [WhiteheadRussell] p. 101.  It says that something is
     either true or not true; there are no in-between values of truth.  This
is
     an essential distinction of our classical logic and is not a theorem of
     intuitionistic logic.  (Contributed by NM, 5-Aug-1993.)
\begin{align}
\vdash ( {\color{formulacolor}\varphi} \vee \lnot {\color{formulacolor}\varphi}
    )\label{eq:exmid}\tag{exmid}
\end{align}
\end{lemma}


\begin{lemma}\otherTitle{falim} $\bot$ implies anything.  (Contributed by FL,
20-Mar-2011.)  (Proof
     shortened by Anthony Hart, 1-Aug-2011.)
\begin{align}
\vdash ( \bot \rightarrow {\color{formulacolor}\varphi} )\label{eq:falim}\tag{falim}
\end{align}
\end{lemma}


\begin{lemma}\otherTitle{idd} Principle of identity with antecedent. 
(Contributed by NM,
     26-Nov-1995.)
\begin{align}
\vdash ( {\color{formulacolor}\varphi} \rightarrow ( {\color{formulacolor}\psi}
    \rightarrow {\color{formulacolor}\psi} ) )\label{eq:idd}\tag{idd}
\end{align}
\end{lemma}


\begin{lemma}\otherTitle{imim2i} Inference adding common antecedents in an
implication.  (Contributed by
       NM, 5-Aug-1993.)
\begin{align}
\vdash ( {\color{formulacolor}\varphi} \rightarrow {\color{formulacolor}\psi}
    )\quad\Rightarrow\quad \vdash ( ( {\color{formulacolor}\chi} \rightarrow
    {\color{formulacolor}\varphi} ) \rightarrow ( {\color{formulacolor}\chi}
    \rightarrow {\color{formulacolor}\psi} ) )\label{eq:imim2i}\tag{imim2i}
\end{align}
\end{lemma}


\begin{lemma}\otherTitle{imim2i} Inference adding common antecedents in an
implication.  (Contributed by
       NM, 5-Aug-1993.)
\begin{align}
\vdash ( {\color{formulacolor}\varphi} \rightarrow {\color{formulacolor}\psi}
    )\quad\Rightarrow\quad \vdash ( ( {\color{formulacolor}\chi} \rightarrow
    {\color{formulacolor}\varphi} ) \rightarrow ( {\color{formulacolor}\chi}
    \rightarrow {\color{formulacolor}\psi} ) )\label{eq:imim2i}\tag{imim2i}
\end{align}
\end{lemma}


\begin{lemma}\otherTitle{impbii} Infer an equivalence from an implication and
its converse.  (Contributed
       by NM, 5-Aug-1993.)
\begin{align}
\vdash ( {\color{formulacolor}\varphi} \rightarrow {\color{formulacolor}\psi}
    )\quad\&\quad \vdash ( {\color{formulacolor}\psi} \rightarrow
    {\color{formulacolor}\varphi} )\quad\Rightarrow\quad \vdash (
    {\color{formulacolor}\varphi} \leftrightarrow {\color{formulacolor}\psi}
    )\label{eq:impbii}\tag{impbii}
\end{align}
\end{lemma}


\begin{lemma}\otherTitle{leid} 'Less than or equal to' is reflexive. 
(Contributed by NM,
     18-Aug-1999.)
\begin{align}
\vdash ( {\color{classcolor}A} \in \mathbb{R} \rightarrow {\color{classcolor}A}
    \le {\color{classcolor}A} )\label{eq:leid}\tag{leid}
\end{align}
\end{lemma}


\begin{lemma}\otherTitle{notbii} Negate both sides of a logical equivalence. 
(Contributed by NM,
       5-Aug-1993.)  (Proof shortened by Wolf Lammen, 19-May-2013.)
\begin{align}
\vdash ( {\color{formulacolor}\varphi} \leftrightarrow
    {\color{formulacolor}\psi} )\quad\Rightarrow\quad \vdash ( \lnot
    {\color{formulacolor}\varphi} \leftrightarrow \lnot
    {\color{formulacolor}\psi} )\label{eq:notbii}\tag{notbii}
\end{align}
\end{lemma}


\begin{lemma}\otherTitle{notnot} Double negation.  Theorem *4.13 of
[WhiteheadRussell] p. 117.
     (Contributed by NM, 5-Aug-1993.)
\begin{align}
\vdash ( {\color{formulacolor}\varphi} \leftrightarrow \lnot \lnot
    {\color{formulacolor}\varphi} )\label{eq:notnot}\tag{notnot}
\end{align}
\end{lemma}


\begin{lemma}\otherTitle{olc} Introduction of a disjunct.  Axiom *1.3 of
[WhiteheadRussell] p. 96.
     (Contributed by NM, 30-Aug-1993.)
\begin{align}
\vdash ( {\color{formulacolor}\varphi} \rightarrow ( {\color{formulacolor}\psi}
    \vee {\color{formulacolor}\varphi} ) )\label{eq:olc}\tag{olc}
\end{align}
\end{lemma}


\begin{lemma}\otherTitle{orbi12i} Infer the disjunction of two equivalences. 
(Contributed by NM,
       5-Aug-1993.)
\begin{align}
\vdash ( {\color{formulacolor}\varphi} \leftrightarrow
    {\color{formulacolor}\psi} )\quad\&\quad \vdash (
    {\color{formulacolor}\chi} \leftrightarrow {\color{formulacolor}\theta}
    )\quad\Rightarrow\quad \vdash ( ( {\color{formulacolor}\varphi} \vee
    {\color{formulacolor}\chi} ) \leftrightarrow ( {\color{formulacolor}\psi}
    \vee {\color{formulacolor}\theta} ) )\label{eq:orbi12i}\tag{orbi12i}
\end{align}
\end{lemma}


\begin{lemma}\otherTitle{orbi1i} Inference adding a right disjunct to both
sides of a logical
       equivalence.  (Contributed by NM, 5-Aug-1993.)
\begin{align}
\vdash ( {\color{formulacolor}\varphi} \leftrightarrow
    {\color{formulacolor}\psi} )\quad\Rightarrow\quad \vdash ( (
    {\color{formulacolor}\varphi} \vee {\color{formulacolor}\chi} )
    \leftrightarrow ( {\color{formulacolor}\psi} \vee
    {\color{formulacolor}\chi} ) )\label{eq:orbi1i}\tag{orbi1i}
\end{align}
\end{lemma}


\begin{lemma}\otherTitle{orbi2i} Inference adding a left disjunct to both sides
of a logical
       equivalence.  (Contributed by NM, 5-Aug-1993.)  (Proof shortened by
Wolf
       Lammen, 12-Dec-2012.)
\begin{align}
\vdash ( {\color{formulacolor}\varphi} \leftrightarrow
    {\color{formulacolor}\psi} )\quad\Rightarrow\quad \vdash ( (
    {\color{formulacolor}\chi} \vee {\color{formulacolor}\varphi} )
    \leftrightarrow ( {\color{formulacolor}\chi} \vee
    {\color{formulacolor}\psi} ) )\label{eq:orbi2i}\tag{orbi2i}
\end{align}
\end{lemma}


\begin{lemma}\otherTitle{orc} Introduction of a disjunct.  Theorem *2.2 of
[WhiteheadRussell] p. 104.
     (Contributed by NM, 30-Aug-1993.)
\begin{align}
\vdash ( {\color{formulacolor}\varphi} \rightarrow (
    {\color{formulacolor}\varphi} \vee {\color{formulacolor}\psi} )
    )\label{eq:orc}\tag{orc}
\end{align}
\end{lemma}


\begin{lemma}\otherTitle{orcom} Commutative law for disjunction.  Theorem *4.31
of [WhiteheadRussell]
     p. 118.  (Contributed by NM, 5-Aug-1993.)  (Proof shortened by Wolf
     Lammen, 15-Nov-2012.)
\begin{align}
\vdash ( ( {\color{formulacolor}\varphi} \vee {\color{formulacolor}\psi} )
    \leftrightarrow ( {\color{formulacolor}\psi} \vee
    {\color{formulacolor}\varphi} ) )\label{eq:orcom}\tag{orcom}
\end{align}
\end{lemma}


\begin{lemma}\otherTitle{ordir} Distributive law for disjunction.  (Contributed
by NM, 12-Aug-1994.)
\begin{align}
\vdash ( ( ( {\color{formulacolor}\varphi} \wedge {\color{formulacolor}\psi} )
    \vee {\color{formulacolor}\chi} ) \leftrightarrow ( (
    {\color{formulacolor}\varphi} \vee {\color{formulacolor}\chi} ) \wedge (
    {\color{formulacolor}\psi} \vee {\color{formulacolor}\chi} ) )
    )\label{eq:ordir}\tag{ordir}
\end{align}
\end{lemma}


\begin{lemma}\otherTitle{oridm} Idempotent law for disjunction.  Theorem *4.25
of [WhiteheadRussell]
     p. 117.  (Contributed by NM, 5-Aug-1993.)  (Proof shortened by Andrew
     Salmon, 16-Apr-2011.)  (Proof shortened by Wolf Lammen, 10-Mar-2013.)
\begin{align}
\vdash ( ( {\color{formulacolor}\varphi} \vee {\color{formulacolor}\varphi} )
    \leftrightarrow {\color{formulacolor}\varphi} )\label{eq:oridm}\tag{oridm}
\end{align}
\end{lemma}


\begin{lemma}\otherTitle{pm2.61} Theorem *2.61 of [WhiteheadRussell] p. 107. 
Useful for eliminating an
     antecedent.  (Contributed by NM, 5-Aug-1993.)  (Proof shortened by Wolf
     Lammen, 22-Sep-2013.)
\begin{align}
\vdash ( ( {\color{formulacolor}\varphi} \rightarrow {\color{formulacolor}\psi}
    ) \rightarrow ( ( \lnot {\color{formulacolor}\varphi} \rightarrow
    {\color{formulacolor}\psi} ) \rightarrow {\color{formulacolor}\psi} )
    )\label{eq:pm2.61}\tag{pm2.61}
\end{align}
\end{lemma}


\begin{lemma}\otherTitle{pm2.86i} Inference based on {\tt pm2.86}. 
(Contributed by NM, 5-Aug-1993.)  (Proof
       shortened by Wolf Lammen, 3-Apr-2013.)
\begin{align}
\vdash ( ( {\color{formulacolor}\varphi} \rightarrow {\color{formulacolor}\psi}
    ) \rightarrow ( {\color{formulacolor}\varphi} \rightarrow
    {\color{formulacolor}\chi} ) )\quad\Rightarrow\quad \vdash (
    {\color{formulacolor}\varphi} \rightarrow ( {\color{formulacolor}\psi}
    \rightarrow {\color{formulacolor}\chi} ) )\label{eq:pm2.86i}\tag{pm2.86i}
\end{align}
\end{lemma}


\begin{lemma}\otherTitle{pm4.44} Theorem *4.44 of [WhiteheadRussell] p. 119. 
(Contributed by NM,
     3-Jan-2005.)
\begin{align}
\vdash ( {\color{formulacolor}\varphi} \leftrightarrow (
    {\color{formulacolor}\varphi} \vee ( {\color{formulacolor}\varphi} \wedge
    {\color{formulacolor}\psi} ) ) )\label{eq:pm4.44}\tag{pm4.44}
\end{align}
\end{lemma}


\begin{lemma}\otherTitle{pm4.45} Theorem *4.45 of [WhiteheadRussell] p. 119. 
(Contributed by NM,
     3-Jan-2005.)
\begin{align}
\vdash ( {\color{formulacolor}\varphi} \leftrightarrow (
    {\color{formulacolor}\varphi} \wedge ( {\color{formulacolor}\varphi} \vee
    {\color{formulacolor}\psi} ) ) )\label{eq:pm4.45}\tag{pm4.45}
\end{align}
\end{lemma}


\begin{lemma}\otherTitle{pm4.45im} Conjunction with implication.  Compare
Theorem *4.45 of [WhiteheadRussell]
     p. 119.  (Contributed by NM, 17-May-1998.)
\begin{align}
\vdash ( {\color{formulacolor}\varphi} \leftrightarrow (
    {\color{formulacolor}\varphi} \wedge ( {\color{formulacolor}\psi}
    \rightarrow {\color{formulacolor}\varphi} ) ) )\label{eq:pm4.45im}\tag{pm4.45im}
\end{align}
\end{lemma}


\begin{lemma}\otherTitle{pm4.56} Theorem *4.56 of [WhiteheadRussell] p. 120. 
(Contributed by NM,
     3-Jan-2005.)
\begin{align}
\vdash ( ( \lnot {\color{formulacolor}\varphi} \wedge \lnot
    {\color{formulacolor}\psi} ) \leftrightarrow \lnot (
    {\color{formulacolor}\varphi} \vee {\color{formulacolor}\psi} )
    )\label{eq:pm4.56}\tag{pm4.56}
\end{align}
\end{lemma}


\begin{lemma}\otherTitle{simpl} Elimination of a conjunct.  Theorem *3.26
(Simp) of [WhiteheadRussell]
     p. 112.  (Contributed by NM, 5-Aug-1993.)  (Proof shortened by Wolf
     Lammen, 13-Nov-2012.)
\begin{align}
\vdash ( ( {\color{formulacolor}\varphi} \wedge {\color{formulacolor}\psi} )
    \rightarrow {\color{formulacolor}\varphi} )\label{eq:simpl}\tag{simpl}
\end{align}
\end{lemma}


\begin{lemma}\otherTitle{syl} An inference version of the transitive laws for
implication {\tt imim2} and
       {\tt imim1}, which Russell and Whitehead call "the principle of the
       syllogism...because...the syllogism in Barbara is derived from them"
       (quote after Theorem *2.06 of [WhiteheadRussell] p. 101).  Some authors
       call this law a "hypothetical syllogism."

       (A bit of trivia: this is the most commonly referenced assertion in our
       database.  In second place is {\tt eqid}, followed by {\tt syl2anc},
       {\tt adantr}, {\tt syl3anc}, and {\tt ax-mp}.  The Metamath program
command 'show
       usage' shows the number of references.)  (Contributed by NM,
       5-Aug-1993.)  (Proof shortened by O'Cat, 20-Oct-2011.)  (Proof
shortened
       by Wolf Lammen, 26-Jul-2012.)
\begin{align}
\vdash ( {\color{formulacolor}\varphi} \rightarrow {\color{formulacolor}\psi}
    )\quad\&\quad \vdash ( {\color{formulacolor}\psi} \rightarrow
    {\color{formulacolor}\chi} )\quad\Rightarrow\quad \vdash (
    {\color{formulacolor}\varphi} \rightarrow {\color{formulacolor}\chi}
    )\label{eq:syl}\tag{syl}
\end{align}
\end{lemma}


\begin{lemma}\otherTitle{sylbi} A mixed syllogism inference from a
biconditional and an implication.
       Useful for substituting an antecedent with a definition.  (Contributed
       by NM, 5-Aug-1993.)
\begin{align}
\vdash ( {\color{formulacolor}\varphi} \leftrightarrow
    {\color{formulacolor}\psi} )\quad\&\quad \vdash (
    {\color{formulacolor}\psi} \rightarrow {\color{formulacolor}\chi}
    )\quad\Rightarrow\quad \vdash ( {\color{formulacolor}\varphi} \rightarrow
    {\color{formulacolor}\chi} )\label{eq:sylbi}\tag{sylbi}
\end{align}
\end{lemma}


\begin{lemma}\otherTitle{sylibr} A mixed syllogism inference from an
implication and a biconditional.
       Useful for substituting a consequent with a definition.  (Contributed
by
       NM, 5-Aug-1993.)
\begin{align}
\vdash ( {\color{formulacolor}\varphi} \rightarrow {\color{formulacolor}\psi}
    )\quad\&\quad \vdash ( {\color{formulacolor}\chi} \leftrightarrow
    {\color{formulacolor}\psi} )\quad\Rightarrow\quad \vdash (
    {\color{formulacolor}\varphi} \rightarrow {\color{formulacolor}\chi}
    )\label{eq:sylibr}\tag{sylibr}
\end{align}
\end{lemma}


\begin{lemma}\otherTitle{trant} A true antecedent can be removed.  (Contributed
by FL, 16-Apr-2012.)
\begin{align}
\vdash ( ( \top \rightarrow {\color{formulacolor}\varphi} ) \leftrightarrow
    {\color{formulacolor}\varphi} )\label{eq:trant}\tag{trant}
\end{align}
\end{lemma}


